\section{Methods}

\subsection{Data sources, patient selection, and study populations}


Anonymized individual-level, population-scale data were accessed through NHS England’s Secure Data Environment (SDE)  service for England (https://digital.nhs.uk/services/secure-data-environment-service).

Details of data sources, patient selection, and study populations are given in the supplementary material. A summary is presented here.

All analysis was performed in SDE using patient level data from seven internally hosted datasets, combining  data from the Sentinel Stroke National Audit Programme (SSNAP), demographic data, Hospital Episode Statistics (HES), Covid infection \& vaccination, and home-time data (days spent alive outside NHS inpatient care). Analysis was performed using Python and the Python Packages pandas \cite{mckinney2011pandas}, matplotlib \cite{barrett2005matplotlib}, statsmodels \cite{seabold2010statsmodels}, XGBoost \cite{XGBoost2026Python}, and SHAP \cite{shap_documentation}.

Three populations were used in the analysis:

\begin{enumerate}
    
    \item The \emph{full stroke population} (n = 388,345) contained patients that had a stroke and attended an acute stroke team between 1\textsuperscript{st} January 2018 and 31\textsuperscript{st} March 2023.
    
    \item The \emph{matched population} of stroke patients (n = 236,630) were a set of stroke patients used to create a population of matched control patients (n = 236,630). Only patients arriving after 1\textsuperscript{st}November 2019 were included, to remove survival bias in the control data set (control patients were only available if they were alive in November 2019).

    \item A \emph{reperfusion analysis cohort} of stroke patients (n = 136,395) eligible for consideration for reperfusion therapy (thrombolysis or thrombectomy). All patients had ischaemic stroke, presented to a stroke team within 6 hours of known onset, and experienced stroke onset outside the hospital.
    
\end{enumerate}

\subsection{Descriptive statistics}

Descriptive statistics were generated, comparing that matched population of stroke and control patients. 


\subsection{Survival analysis}

Survival analysis was performed on the \textit{full stroke population} using a Cox proportional hazard ratio regression model. Features included in the model were age, sex, ethnicity, IMD quintiles, stroke severity, prior disability, stroke type, treatment received, and covid status (infection and/or vaccination).

\subsection{Explainable machine learning for prediction of stroke patient outcome and location in the first year following their stroke}

Explainable machine learning was used to predict time in care, at-home time, and time dead in the \textit{full stroke population} and the \textit{reperfusion analysis cohort} population. Features were selected for the \textit{full stroke population} model as those features identified as improving prediction of outcomes, or of being of primary interest (deprivation, ethnicity, treatment, COVID status; see Supplementary Material for details), whereas features for the \textit{reperfusion analysis} model were treatment features, and identified confounding features (those features associated with outcome directly, \textit{and} with use of reperfusion treatment previously identified \cite{pearn_thrombolysis_2026}).

We selected XGBoost (eXtreme Gradient Boosting) \cite{chen_xgboost_2016} as the primary modelling approach due to its computational efficiency and strong performance on tabular data, often requiring minimal hyperparameter tuning \cite{shwartz-ziv_tabular_2022}. This reduces the risk of overfitting arising from extensive model optimisation. XGBoost is also well suited to capturing complex non-linear relationships and feature interactions. However, XGBoost models are inherently limited in interpretability at both the individual and group levels. To address this, we applied SHapley Additive exPlanations (SHAP) \cite{lundberg_consistent_2019} to quantify the contribution and directional effect of each clinical variable. We fitted three separate XGBoost models to predict time spent in NHS inpatient care, time spent dead, and time spent alive within one year following stroke. For each model, SHAP defines a baseline prediction (common to all individuals), which is adjusted by the sum of feature-specific SHAP values for each individual to yield the final prediction. A positive SHAP value indicates that the corresponding feature increases the predicted duration, whereas a negative value indicates a reduction.

Accuracy of models was tested using 5-fold cross-validation, where the data was split into five 80:20 training:test data sets, with each patient being in one, and only one, test set.  Accuracy is reported as R-squared.